\chapter{Testování a evaluace řešení}
\label{sec:testing}

Testování systému NPC probíhalo manuálně v rámci herní scény. 
Cílem testování bylo ověřit, zda implementovaná chování odpovídá návrhu a zda jednotlivé mechanismy fungují v očekávaných situacích.

\begin{table}[H]
\centering
\begin{tabular}{|p{0.28\textwidth}|p{0.42\textwidth}|p{0.18\textwidth}|}
\hline
\textbf{Scénář} & \textbf{Očekávané chování} & \textbf{Výsledek} \\
\hline
Přechody mezi stavy & NPC přechází mezi stavy \texttt{Passive}, \texttt{Investigating} a \texttt{Attacking}. & Ověřeno \\
\hline
Zraková percepce & Detekce hráče vede k přechodu do stavu \texttt{Attacking}. & Ověřeno \\
\hline
Ztráta vizuálního kontaktu & NPC se po ztrátě cíle vrací do pasivního chování. & Ověřeno \\
\hline
Sluchová percepce & Zvukový podnět vede k přechodu do stavu \texttt{Investigating}. & Ověřeno \\
\hline
Hlídkování NPC & NPC se pohybuje po přiřazené hlídkové trase. & Ověřeno \\
\hline
Změna hlídkové trasy & NPC po splnění podmínek změní aktuální trasu. & Ověřeno \\
\hline
Fázově řízené trasy & NPC upravují výběr tras podle fáze hry. & Ověřeno \\
\hline
Interakce s objekty & NPC dokáže interagovat s objekty (např. otevírání dveří). & Ověřeno \\
\hline
Sebrání předmětu & NPC dokáže sebrat interaktivní předmět (např. klíč). & Ověřeno \\
\hline
Položení předmětu & NPC po sebrání přesune předmět na jiné místo. & Ověřeno \\
\hline
Sdílení informací & NPC si předávají informace o předmětech. & Neověřeno \\
\hline
\end{tabular}
\caption{Přehled manuálně ověřených scénářů}
\label{tab:testovaci-scenare}
\end{table}

Z tabulky~\ref{tab:testovaci-scenare} vyplývá, že základní mechanismy systému fungují dle očekávání. 
Byla ověřena zejména správná funkce percepce, přechodů mezi stavy a hlídkovacího chování NPC.

Naopak některé pokročilejší části systému nebyly v rámci testování spolehlivě ověřeny. 
Konkrétně se jedná o manipulaci s předměty po jejich sebrání a sdílení informací mezi NPC. 
Tyto části proto nelze považovat za plně otestované a představují omezení aktuální implementace.